% ME310 Report Template -- Started 28 July 2006  Mark Cutkosky
% 24Nov2016: Changes in ordering and content to go with the revised
%  schedule for ME310A in 2016-17
%%%%%%%%%%%%%%%%%%%%%%%%%%%%%%%%%%%

% Rename this file to whatever you like (e.g. OurFallDoc.tex) and modify title, authors
% etc. below. If you rename section files (e.g. Ch2context.tex) you'll need to change
% the \include{ } calls in this document as well.

%%%%%%%%%BEGIN DOCUMENT STYLE SETTINGS%%%%%%%%%%%
% Don't modify this stuff unless you know what you're doing...
% We are using the "memoir" class, a widely used set of macros for book-like documents.
% If you get errors that you are missing the "memoir" package you can download and 
% install it:   http://www.ctan.org/tex-archive/macros/latex/contrib/memoir/

% memoir document class for standard USA letter paper, printed one side
\documentclass[11pt,letterpaper,oneside]{memoir}
%% Use the 'draft' setting below if things are bogging down in Overleaf. 
%% It replaces figures with placeholders, to compile much faster.
%\documentclass[11pt,letterpaper,oneside,draft]{memoir}
\chapterstyle{section}
%\pagestyle{companion}    % If you want to experiment with fancier page headers 
\usepackage{graphicx}        % standard LaTeX graphics 
\usepackage{color}               % support for colored fonts
\usepackage{url}  \urlstyle{same}     % deal with url strings in bibliography

%Optional fonts if you don't like the default in Latex
%\usepackage[T1]{fontenc}
\usepackage{mathptmx}  %This will give you Times Roman, like default in MS Word.
%\usepackage{charter}  %This will give you a slightly bolder Charter font 

%More special packages to help deal with long requirements tables 
%that might span multiple pages.
\usepackage{multirow} %deal with merged cells in tables
\usepackage{supertabular}
\usepackage{longtable}
\usepackage{morefloats}

\usepackage[pdftex,           %Hyperlink cross references, URLs, etc.
    pdfsubject={ME310 Documentation},
    colorlinks={true},
    linkcolor={black},
    citecolor={blue},
    bookmarksopenlevel=1,
]{hyperref}				  

%The file "me310.sty" should be in the same directory as this file.
% It contains formatting for page setup, titlepage, glossary, references, etc.
\usepackage{me310}           
%%%%%%%END DOCUMENT STYLE SETTINGS%%%%%%%%%%


%%%%%%%%%%BEGIN TITLE PAGE%%%%%%%%%%%%%%%%
%% Some teams will want to replace this rather plain cover sheet with a fancier
%% one that they create by hand in Powerpoint or similar program.

%Replace the strings below with what's right  for you.
%%Insert your Document Title here. Use \\ to force a newline.
\title{An Internet Appliance\\
for Individual Well Being}

\team{Our Team Name} 		% Insert your Team Project Name here.

%% Insert Fall, Winter, Spring here:
\quarter{Fall Design Document}

%%Enter local + global team members' names here:
\author{  
    Person One, Person Two, Person Three\\
    Person Four, Person Five, Person Six, Person Seven\\
    } %end authors

%% If you don't want it to use the printing date, replace "\today"
%% with the date that you want.
\date{\today}
%%%%%%%%%%%END TITLE PAGE%%%%%%%%%%%%%


%%%%%%%%%BEGIN ANY CUSTOM ABBREVIATIONS%%%%%%
% Define any abbreviations that will apply throughout the document
% to save typing. Examples:
\def\pmt{{\em Papier M\^{a}ch\`{e}}}  %Define "\pmt" to print "Papier Mache" with accents +1space
\def\cbike{{\em Casterbike} \,}  %Define "\cbike" to print "Casterbike " italicized +1space
%%%%%%%%%END CUSTOM ABBREVIATIONS%%%%%%%%%

%%%%%%%%COMMENTS%%%%%%%%%%%%%%%%
%Allow comments ('remarks') to be shown or hidden.
% Put optional text in \begin{remark} ... \end{remark} environments.
%The usage is counterintuitive: \commentsoff makes them visible; \commentson hides them.
%\commentson{remark} %Don't print remarks.
\commentsoff{remark}  %Do print remarks!


%%%%%%%%%%%%%%%%%%%%%%%%%%%%%%%%%%%%%%%%
%   BEGIN THE MAIN DOCUMENT
%%%%%%%%%%%%%%%%%%%%%%%%%%%%%%%%%%%%%%%%
\begin{document}

%If you want a figure on the cover page, this is where it goes.
%8 cm is about max figure height before messing up title spacing.
%If making your own fancy coverpage (e.g. in Photoshop) then comment out 
%the \includegraphics{ } and use the \vspace{ } command.
\begin{figure}[t]
\centering
  %An example cover image, from 2008 ME310 Kodak project
  \includegraphics[height= 8cm]{Figures/Ch1/KodakThing.pdf}
%\vspace{6 cm}    %Use this instead if you have no cover picture 
\end{figure}

%%%%%%%%%%%%%%%%%%%%%%%%%%%%%%%%%%%%%%
%Make the title page using arguments defined above.
\titlep

%%%%%%%%%%%%%%%%%%%%%%%%
% Load file "Executive.tex" for the Executive summary.
% Remember, this is a stand-alone section for executives to read.
\include{Executive}

%%%%%%%%%%%%%%%%%%%%%%%%
% TOC and LOF are automatically generated. Depending on Latex environment,
% you might have to "compile" THREE times to update all .aux files, TOC, 
% cross-references etc. 
% This is not an issue with Overleaf, which continuously recompiles.


% Make Table of Contents title smaller than a normal Chapter heading:
\renewcommand{\chaptitlefont}{\normalfont\Large\bfseries}
\newpage
\tableofcontents*  %asterisk to prevent it from getting a number

% Optional Lists of Figures and Tables:
\newpage
\listoffigures*  %For this you probably want to add the [short-headings] to captions.
\listoftables*

%%%%%% Example of an optionally printed "remark"
\begin{remark}
\color{blue}
When creating figures and tables, use "caption[short caption for ToC]\{Long and descriptive caption, as one finds in Nature or Science, etc.\}" to get a nice formatting effect.
\end{remark}


%Back to normal size for subsequent sections
\renewcommand{\chaptitlefont}{\normalfont\Huge\bfseries}
%%%%%%%%%%%%%%%%%%%%%%%%

% Set up the Glossary. The template is looking for a file called
% "glossaryterms.tex" with glossary terms and definitions.
\newpage
\section*{Glossary}
\addcontentsline{toc}{section}{Glossary}
\label{sec-glossary}
\begin{description}
\input{glossaryterms}   % input the list "glossaryterms.tex"
\end{description}

%%%%%% Example of an optionally printed "remark"
\begin{remark}
\color{blue}
It's a sign of a successful team that the glossary becomes extensive. Define any non-obvious or invented terms. For example, if you reference something by an acronym, that might be a glossary term. Teams also coin terms to describe design features. Don't define obvious stuff (axle, keyboard).  
\normalcolor
\end{remark}


%%%%%%%%%%%%%%%%%%%%%%%%%%%%%%%%
%% On to the main sections....  Just comment out the \input{} line
%% for any chapters that aren't ready yet.

%%%%%%%%%%%%%%%%%%%%%%%%%%%%%%%%
\include{Context}        %Project and corporate background

%%%%%%%%%%%%%%%%%%%%%%%%%%%%%%%%
\include{Users}          %Users, need-finding, personas

%%%%%%%%%%%%%%%%%%%%%%%%%%%%%%%%
\include{Benchmarking}   %Relevant technology, systems

%%%%%%%%%%%%%%%%%%%%%%%%%%%%%%%%
\include{Prototyping}    %Early Experience and Functional prototypes

%%%%%%%%%%%%%%%%%%%%%%%%%%%%%%%%
\include{Vision}         %Current vision or proposal for the design

%%%%%%%%%%%%%%%%%%%%%%%%%%%%%%%%
\include{Requirements}   %Given and Identified requirements for vision

%%%%%%%%%%%%%%%%%%%%%%%%%%%%%%%%
\include{Reflections}    %Current and future plans.

%%%%%%%%%%%%%%%%%%%%%%%%%%%%%%%%
\include{Resources}  

%%%%%BEGIN BIBLIOGRAPHY%%%%%%%%%
% This is a two-step process in which you first create a ".bib" file, which is processed
% by bibtex into a ".bbl" file for loading into the document. 
% Many online journals and databases now have a feature to automatically download Bib
% citations.  BibDesk is a handy (free) program for Macintosh for managing them.
% EndNote and Refworks (is free to Stanford students) are good alternatives.
\bibliographystyle{plainurl310}     %Modified slightly from plainurl.bst 
\bibliography{me310reportfall}  % Look for file called "me310report.bib" 
%
% Alternatively, you can create your own bibliography list by hand.
% In that case, comment out the last line above and replace with  \input{OurBibFile} 
%%%%%END BIBLIOGRAPHY%%%%%%%%%

%%%%%BEGIN APPENDIX SECTIONS%%%%%%%%%%%%%%%
\begin{appendices}
 \include{Appendix1}       % There could be multiple appendix files like this
 \end{appendices}
 %%%%%END APPENDIX SECTIONS%%%%%%%%%%%%%%%
 
\end{document}

